\section*{Bab \ref{ch:g2:problems} --- Menyelesaikan Soal Cerita}

\begin{solution}{exo:g2:problems:1}
Digabungkan: penjumlahan. $50$ permen, $12$ dimakan: pengurangan.
$4$ kotak berisi $5$: perkalian. Berapa lebihnya: pengurangan.
\end{solution}

\begin{solution}{exo:g2:problems:2}
$156 - 48 = 108$. Yang tinggal $108$ mobil.
\end{solution}

\begin{solution}{exo:g2:problems:3}
$5 \times 4 = 20$ bunga.
\end{solution}

\begin{solution}{exo:g2:problems:4}
Batang pembanding: Tono $45$, batang Mia lebih pendek $17$:
$45 - 17 = 28$ kartu.
\end{solution}

\begin{solution}{exo:g2:problems:5}
Botol yang diterima: $4 \times 6 = 24$; yang tersisa:
$24 - 9 = 15$ botol.
\end{solution}

\begin{solution}{exo:g2:problems:6}
Bilangan yang tidak berguna adalah umurnya ($8$). Kelereng:
$23 + 15 = 38$.
\end{solution}

\begin{solution}{exo:g2:problems:7}
Pai: $7 + 10 = 17$. Hari Sabtu, semua kue:
$12 + 7 = 19$.
\end{solution}

\begin{solution}{exo:g2:problems:8}
Merah $8$, hijau $4$, kuning $6$: seluruhnya $8 + 4 + 6 = 18$ buah.
Merah yang paling banyak.
\end{solution}

\begin{solution}{exo:g2:problems:9}
$80 - 35 = 45$ cm tersisa.
\end{solution}

\begin{solution}{exo:g2:problems:10}
Pagi: $5 \times 4 = 20$; sore: $3 \times 4 = 12$; \omterm{def:g1:addition:def}{jumlahnya}
$20 + 12 = 32$ kue. (Atau: $5 + 3 = 8$ loyang berisi $4$, lalu
$8 \times 4 = 32$.)

\end{solution}
